% ===== Preâmbulo compartilhado (Tectonic / XeTeX) — documentos da FTL094 =====
\documentclass[11pt,a4paper]{article}
\usepackage{fontspec}
\usepackage[brazil]{babel}
\usepackage{geometry}
\geometry{a4paper,margin=2.5cm}
\usepackage{amsmath,amssymb}
\usepackage{graphicx}
\usepackage{booktabs}
\usepackage{array}
\usepackage{longtable}
\usepackage{tabularx}
\usepackage{float}
\usepackage{enumitem}
\usepackage{xcolor}
\usepackage{tikz}
\usetikzlibrary{arrows.meta,positioning,automata,shapes.geometric,fit,backgrounds,calc}
\usepackage{pgfplots}
\pgfplotsset{compat=1.18}
\usepackage{listings}
\usepackage{caption}
\usepackage{fancyhdr}
\usepackage[hidelinks]{hyperref}

% ----- cores -----
\definecolor{ufamblue}{RGB}{0,51,102}
\definecolor{codegreen}{rgb}{0,0.5,0}
\definecolor{codeblue}{rgb}{0,0,0.7}
\definecolor{codered}{rgb}{0.6,0,0}
\definecolor{lgray}{rgb}{0.95,0.95,0.95}

% ----- listagens C -----
\lstdefinestyle{c}{language=C,basicstyle=\ttfamily\footnotesize,
  keywordstyle=\color{codeblue}\bfseries,commentstyle=\color{codegreen}\itshape,
  stringstyle=\color{codered},numbers=left,numberstyle=\tiny\color{gray},
  frame=single,breaklines=true,tabsize=2,showstringspaces=false,columns=fullflexible}

% ----- constantes do projeto -----
\newcommand{\projnome}{Sistema de Navegação Autônoma para o Robô Móvel Khepera~IV}
\newcommand{\disciplina}{FTL094 --- Engenharia de Software de Tempo Real}
\newcommand{\periodo}{Período Letivo 2026/1}
\newcommand{\versaodoc}{1.0}

% ----- cabeçalho/rodapé -----
\pagestyle{fancy}\fancyhf{}
\rhead{\small \tipodoc}
\lhead{\small Khepera~IV --- FTL094}
\cfoot{\thepage}
\renewcommand{\headrulewidth}{0.4pt}

% ----- coluna X centralizada p/ tabularx -----
\newcolumntype{C}{>{\centering\arraybackslash}X}

% ----- capa padronizada -----
\newcommand{\fazcapa}{%
\begin{titlepage}\centering
{\large\bfseries UNIVERSIDADE FEDERAL DO AMAZONAS (UFAM)}\\[3pt]
{\normalsize Centro de P\&D em Tecnologia Eletrônica e da Informação --- CETELI}\\[2pt]
{\normalsize Programa de Pós-Graduação em Engenharia Elétrica --- PPGEE}\\[26pt]
{\large \disciplina}\\[2pt]
{\normalsize \periodo}\\[64pt]
\rule{\textwidth}{1.2pt}\\[12pt]
{\Huge\bfseries\color{ufamblue} \tipodoc}\\[14pt]
{\Large \projnome}\\[10pt]
\rule{\textwidth}{1.2pt}\\[48pt]
{\large\bfseries Equipe}\\[8pt]
{\normalsize Gabriel Henrique de Souza Nazaré}\\[3pt]
{\normalsize Luan Alexandre Ramos Barreto}\\[3pt]
{\normalsize Matheus Rocha Canto}\\[3pt]
{\normalsize Matheus Serrão Uchôa}\\[3pt]
{\normalsize Ricardo Mendonça Braz}\\[3pt]
{\normalsize Yan Fernandes da Silva}\\[40pt]
{\normalsize Docente responsável: Prof. Dr.~[nome do docente]}\\
\vfill
{\small Documento de tipo \emph{\tipodoc} --- Versão \versaodoc}\\[2pt]
{\small Manaus --- AM, \datadoc}
\end{titlepage}}

% ----- tabela de controle de versões -----
\newcommand{\controleversoes}{%
\section*{Controle de versões}
\begin{tabularx}{\textwidth}{@{}l l l X@{}}
\toprule
\textbf{Versão} & \textbf{Data} & \textbf{Autor} & \textbf{Descrição} \\
\midrule
1.0 & \datadoc & Equipe FTL094 & Emissão inicial do documento. \\
\bottomrule
\end{tabularx}
\bigskip}

\newcommand{\tipodoc}{Código-Fonte}
\newcommand{\datadoc}{9 de julho de 2026}

% estilos de listagem adicionais (mundo e Makefile)
\lstdefinestyle{plain}{basicstyle=\ttfamily\scriptsize,frame=single,breaklines=true,
  numbers=left,numberstyle=\tiny\color{gray},columns=fullflexible,showstringspaces=false}
\lstdefinestyle{mk}{language=make,basicstyle=\ttfamily\footnotesize,frame=single,
  breaklines=true,numbers=left,numberstyle=\tiny\color{gray},columns=fullflexible}

\begin{document}
\fazcapa
\controleversoes
\tableofcontents
\newpage

% =====================================================================
\section{Introdução}
% =====================================================================
\subsection{Objetivo}
Este documento reúne, na íntegra, o \textbf{código-fonte desenvolvido} para o \emph{\projnome}, conforme exigido na entrega final do projeto. As listagens são extraídas diretamente dos arquivos do repositório, garantindo fidelidade ao código efetivamente executado e avaliado no Protocolo de Testes.

\subsection{Escopo}
O software compreende três artefatos: (i) o \textbf{mundo de simulação} (cenário do Webots), (ii) o \textbf{controlador de navegação} em linguagem C e (iii) o \textbf{script de compilação} (\emph{Makefile}). A lógica de navegação --- arquitetura híbrida com máquina de estados --- está descrita no documento de Arquitetura; aqui o foco é o código.

% =====================================================================
\section{Organização e estrutura do código}
% =====================================================================
A estrutura de diretórios segue a convenção de projetos do Webots (pastas \texttt{worlds/} e \texttt{controllers/}):

\begin{lstlisting}[style=plain,numbers=none]
eng_soft_rt/
  worlds/
    first.wbt                 <- mundo de simulacao (cenario)
  controllers/
    goto_avoid/
      goto_avoid.c            <- controlador de navegacao (FSM)
      Makefile                <- script de compilacao
\end{lstlisting}

\begin{table}[H]\centering
\caption{Artefatos de código e métricas.}
\renewcommand{\arraystretch}{1.2}
\begin{tabularx}{\textwidth}{@{}l l c X@{}}
\toprule
\textbf{Arquivo} & \textbf{Linguagem} & \textbf{Linhas} & \textbf{Descrição} \\
\midrule
\texttt{worlds/first.wbt} & VRML/PROTO (Webots) & 83 & Arena, obstáculos, alvo e robô (com GPS+IMU). \\
\texttt{controllers/goto\_avoid/goto\_avoid.c} & C & 163 & Controlador de navegação com máquina de estados. \\
\texttt{controllers/goto\_avoid/Makefile} & Make & 6 & Compilação via cadeia de ferramentas do Webots. \\
\midrule
\textbf{Total} & --- & \textbf{252} & --- \\
\bottomrule
\end{tabularx}
\end{table}

% =====================================================================
\section{Compilação e execução}
% =====================================================================
O controlador em C é compilado pela \textbf{cadeia de ferramentas embarcada no Webots} (MinGW-w64), sem dependências externas. Há duas formas:

\noindent\textbf{(a) Pela interface gráfica:} abrir \texttt{worlds/first.wbt} no Webots e pressionar \emph{Play}; o controlador é compilado automaticamente e a simulação inicia.

\noindent\textbf{(b) Por linha de comando (execução \emph{headless}, usada nos testes):}
\begin{lstlisting}[style=plain,numbers=none]
set WEBOTS_HOME=C:\Program Files\Webots
make -C controllers\goto_avoid                 # compila goto_avoid.exe
webots --batch --mode=fast --no-rendering worlds\first.wbt
\end{lstlisting}
A telemetria periódica é gravada em \texttt{controllers/goto\_avoid/run.log}, permitindo a verificação objetiva dos casos de teste.

% =====================================================================
\section{Listagens}
% =====================================================================
\subsection{Mundo de simulação --- \texttt{worlds/first.wbt}}
Define o sistema de coordenadas, a arena $3\times3$~m, os cinco obstáculos, o marcador de alvo e a instância do robô Khepera~IV, equipada com GPS e unidade inercial no \emph{turret slot} e associada ao controlador \texttt{goto\_avoid}.

\lstinputlisting[style=plain]{../worlds/first.wbt}

\subsection{Controlador de navegação --- \texttt{controllers/goto\_avoid/goto\_avoid.c}}
Implementa o laço de controle periódico e a máquina de estados (\textsc{Navegar}/\textsc{Girar}/\textsc{Contornar}), incluindo percepção (GPS, IMU, IR), leis de controle, cinemática diferencial e telemetria.

\lstinputlisting[style=c]{../controllers/goto_avoid/goto_avoid.c}

\subsection{Script de compilação --- \texttt{controllers/goto\_avoid/Makefile}}
\lstinputlisting[style=mk]{../controllers/goto_avoid/Makefile}

% =====================================================================
\section{Correspondência código $\leftrightarrow$ arquitetura}
% =====================================================================
A Tabela~\ref{tab:map} relaciona os principais trechos do controlador aos componentes da arquitetura e aos requisitos atendidos, facilitando a leitura e a avaliação.

\begin{table}[H]\centering
\caption{Mapa do código para componentes e requisitos.}
\label{tab:map}
\renewcommand{\arraystretch}{1.2}
\begin{tabularx}{\textwidth}{@{}X l l@{}}
\toprule
\textbf{Trecho do código} & \textbf{Componente} & \textbf{Requisitos} \\
\midrule
\texttt{wb\_gps\_get\_values}, \texttt{wb\_inertial\_unit\_*} & Percepção/Localização & RF-04 \\
Leitura dos IR; \texttt{block\_l}, \texttt{block\_r} & Percepção de obstáculos & RF-02 \\
\texttt{switch(state)} (Navegar/Girar/Contornar) & Arbitragem (FSM) & RF-02, RF-06 \\
\texttt{steer = K\_HEADING*err + K\_AVOID*\dots} & Controle (lei de rumo) & RF-01, RF-05 \\
\texttt{vl}, \texttt{vr}, \texttt{wb\_motor\_set\_velocity} & Atuação (cinemática diferencial) & RF-05 \\
\texttt{if (dist < GOAL\_TOL)} & Detecção de chegada/parada & RF-03 \\
\texttt{fprintf(logf, \dots)}, \texttt{fflush} & Telemetria & RF-07 \\
\texttt{while (wb\_robot\_step(dt) != -1)} & Laço periódico (tempo real) & RT-01 \\
\bottomrule
\end{tabularx}
\end{table}

% =====================================================================
\section{Considerações finais}
% =====================================================================
O código apresentado é o artefato verificado no Protocolo de Testes, no qual o sistema foi \textbf{aprovado} (missão concluída sem colisões, com escape de mínimo local e respeito ao período do laço). Os parâmetros de controle estão centralizados como constantes no início do controlador, favorecendo a manutenibilidade (RNF-05) e a futura portabilidade para o robô físico (RNF-04), bastando substituir a camada de localização (GPS$\to$odometria).

\end{document}
